\section{Onde aprofundar}

Se for ler um só livro, comece pelo de Kent Beck (TDD) ou pelo
\emph{Accelerate} (visão de entrega, aula 2). Os títulos marcados com
\pt têm edição em português.

\subsection*{Livros}
\small
\begin{longtable}{@{}p{1.5cm} p{4.4cm} p{2.6cm} p{4cm}@{}}
\toprule
Tema & Livro & Autor(es), ano & Por que ler \\
\midrule \endhead
TDD & \citetitle{livro:beck2002} \pt{} \emph{TDD: Desenvolvimento Guiado por
  Testes} & \citeauthor{livro:beck2002}, \citeyear{livro:beck2002}
  & A origem do fake it, da triangulação e da implementação óbvia \\
TDD & \citetitle{livro:martin2021} \pt{} \emph{Craftsmanship Limpo}
  & \citeauthor{livro:martin2021},
  \citeyear{livro:martin2021} & TDD passo a passo com katas em Java e JUnit \\
TDD & \citetitle{livro:martin2008}, cap.~9 \pt{} \emph{Código Limpo}
  & \citeauthor{livro:martin2008}, \citeyear{livro:martin2008}
  & As três leis e o F.I.R.S.T. \\
TDD & \citetitle{livro:martin2011}, cap.~5 \pt{} \emph{O Codificador Limpo}
  & \citeauthor{livro:martin2011}, \citeyear{livro:martin2011}
  & TDD como disciplina profissional \\
TDD & \citetitle{livro:freeman2009} & \citeauthor{livro:freeman2009},
  \citeyear{livro:freeman2009} & TDD de fora para dentro, com mocks nas
  bordas \\
Testes & \citetitle{livro:aniche2015} \pt & \citeauthor{livro:aniche2015},
  \citeyear{livro:aniche2015} (Casa do Código) & Introdução prática em
  português, com Java \\
Testes & \citetitle{livro:aniche2022} & \citeauthor{livro:aniche2022},
  \citeyear{livro:aniche2022} & Testes por especificação, bordas,
  propriedades e mutação, em Java \\
Testes & \citetitle{livro:langr} & \citeauthor{livro:langr}
  & O Right-BICEP da aula \\
Testes & \citetitle{livro:khorikov2020} & \citeauthor{livro:khorikov2020},
  \citeyear{livro:khorikov2020} & Quando usar mocks, estado vs
  comportamento \\
Testes & \citetitle{livro:meszaros2007} & \citeauthor{livro:meszaros2007},
  \citeyear{livro:meszaros2007} & De onde vem o vocabulário dos dublês de
  teste \\
Testes & \citetitle{livro:feathers2004} \pt{} \emph{Trabalho Eficaz com
  Código Legado} & \citeauthor{livro:feathers2004},
  \citeyear{livro:feathers2004} & Como colocar testes em código que não tem \\
Produto & \citetitle{livro:adzic2011} & \citeauthor{livro:adzic2011},
  \citeyear{livro:adzic2011} & Do pedido de produto a exemplos que viram
  testes \\
\bottomrule
\end{longtable}
\normalsize

\subsection*{Blog posts e artigos}
Todos em inglês e gratuitos. Os dois posts do Uncle Bob são a base das
aulas.

\small
\begin{longtable}{@{}p{1.5cm} p{6.5cm} p{4.5cm}@{}}
\toprule
Tema & Artigo & Autor, ano \\
\midrule \endhead
TDD & \citetitle{artigo:cycles-tdd} & \citeauthor{artigo:cycles-tdd},
  \citeyear{artigo:cycles-tdd} \\
TDD & \citetitle{artigo:tpp} & \citeauthor{artigo:tpp},
  \citeyear{artigo:tpp} \\
Dublês & \citetitle{artigo:mocks-arent-stubs} &
  \citeauthor{artigo:mocks-arent-stubs} \\
Dublês & \citetitle{artigo:testdouble} & \citeauthor{artigo:testdouble} \\
\bottomrule
\end{longtable}
\normalsize

\subsection*{Vídeos}
Palestras e tutoriais, a maioria em inglês; os marcados com \pt são em
português. Se tiver tempo para um só, assista ao do Ian Cooper
(\citetitle{video:cooper-tdd-wrong}).

\small
\begin{longtable}{@{}p{1.5cm} p{7cm} p{4cm}@{}}
\toprule
Tema & Vídeo & Quem \\
\midrule \endhead
TDD & \citetitle{video:cooper-tdd-wrong} (Build Stuff 2013) &
  \citeauthor{video:cooper-tdd-wrong} \\
TDD & \citetitle{video:cooper-tdd-revisited} (NDC Porto 2023) &
  \citeauthor{video:cooper-tdd-revisited} \\
TDD & \citetitle{video:martin-clean-code-4} (aula sobre TDD) &
  \citeauthor{video:martin-clean-code-4} \\
TDD & \citetitle{video:beck-theme-variations} &
  \citeauthor{video:beck-theme-variations} \\
TDD & \citetitle{video:tdd-is-dead-2} & Kent Beck, Martin Fowler e David
  Heinemeier Hansson \\
Testes & \citetitle{video:aniche-escrever-testes} \pt &
  \citeauthor{video:aniche-escrever-testes} \\
\bottomrule
\end{longtable}
\normalsize

\subsection*{Ferramentas e documentação}
Nesta aula: \citetitle{doc:junit5}, \citetitle{doc:assertj},
\citetitle{doc:mockito} (testes unitários), \citetitle{doc:jqwik}
(propriedades) e \citetitle{doc:pitest} (mutação). A lista completa, com
links, está no apêndice de ferramentas.
